\scanpage{12}%
tal de $X$ s'identifie à un sous-groupe de $\deleted{\unsure{\Gamma}}\Gamma_{p,q}$. Si $\unsure{\mathscr C|C|\mathcal C}$ est régulière \inserted{épinglée}, $G$ le grpe de ses automorphismes, alors le grpe fondamental mixte de $(X,G)$ s'identifie à $\Gamma_{p,q}$

\textcircled{5} [Appelons \deleted{\unsure{plan métrique espace|plan métrique}} \textquotedblleft surface \inserted{(riemannienne)} standart\textquotedblright\ une surface \deleted{\unsure{diff}} \inserted{$E$} analytique réelle munie d'une métrique \inserted{riemannienne $>0$} telle que

a) $E$ est $\ne\varnothing$, connexe, simplement connexe

b) Le groupe des isométries de $E$ est transitif sur $E$.

\deleted{\unsure{c)}} On sait alors que, à un facteur constant près sur la métrique, $E$ est isomorphe à un des trois espaces :

1\textsuperscript{o}) La sphère $S^2$ comme espace homogène de $\mathrm{SO}(3,\mathbb R)$
% slanted
\margin{espace sphérique, déplacements sphériques}

2\textsuperscript{o}) Le plan \deleted{euclidien} \inserted{complexe} \inserted{avec $\unsure{\{z\bar z|qz\bar z|\{z,\bar z\}}$} comme espace homogène sous le groupe $\unsure{\mathbb U|U|\mathcal U}.\unsure{C|\mathbb C}$ (1/2 direct) le groupe des déplacements euclidiens
% slanted
\margin{plan (ou espace) euclidien,}

3\textsuperscript{o}) Le demi-plan de Poincaré, \deleted{\unsure{comme|comme espace}} espace homogène de $\mathrm{SL}(2,\mathbb R)/\{\pm1\}$.
% slanted
\margin{plan hyperbolique, dépl. hyperboliques}

Les groupes écrits sont les \uline{comp. connexes} du groupe des isométries -- ou \unsure{encore|ens|mieux}, \inserted{\unsure{formés|formé|fermées}} des isométries qui \uline{conservent} \uline{l'orientation} -- ou encore \unsure{le|la|les}
% archive sheet number: 11

